La derivación de \(G\) lleva esa reciprocidad a una condición geométrica mucho más concreta.\par

La Mecánica Dimensional determina primero qué tipo de objeto puede ser \(G\). Fija su recta de magnitud y su homogeneidad dimensional. El cierre cronológico selecciona después su valor: la dimensión determina la clase de magnitud y el cierre determina su coeficiente adimensional.\par

La selección procede del periodicidad temporal de (108) fases.\par

El ciclo \(108\) determina una frecuencia, un radio de fase y una energía. A partir de esa energía aparece una masa cronológica. Y al construir la longitud de Compton reducida de esa masa se obtiene exactamente el mismo radio del ciclo:\par

\[
\bar\lambda_C=R_{108}.
\]

La fase y la materia ya están diciendo la misma longitud antes de introducir \(G\).\par

Entonces entra la gravedad. Se construye el radio gravitatorio reducido\par

\[
r_g=\frac{Gm}{c^2},
\]

y se exige que el cierre gravitatorio identifique el radio de fase, la longitud de Compton y el radio gravitatorio como una sola escala estructural:\par

\[
r_g=R_{108}=\bar\lambda_C.
\]

Esa condición selecciona unívocamente\par

\[
\theta_{108}
=
\left(\frac{54}{\pi}\right)^2.
\]

De este modo,\par

\[
R_{108}
=
\bar\lambda_C
=
r_g
=
\ell_P.
\]

La función estructural de \(\theta_{108}\) determina la significación de su decimal terminal: muestra qué está haciendo \(G\).\par

\(G\) aparece como el coeficiente necesario para hacer coincidir las cuatro longitudes publicadas por estos lenguajes:\par

\begin{itemize}[leftmargin=*,itemsep=0.35em]
\item el lenguaje de la fase produce \(R_{108}\);
\item el lenguaje cuántico produce \(\bar\lambda_C\);
\item el lenguaje gravitatorio produce \(r_g\);
\item el lenguaje de Planck produce \(\ell_P\).
\end{itemize}

La constante gravitatoria es el mediador que impide que esos cuatro lectores se separen.\par

Dicho de una manera más física: \(G\) expresa cuánto debe responder la geometría para que un cuanto cuya masa procede del reloj \(108\) cierre sobre su propia fase.\par

Es la relación que hace posible que materia, fase y geometría sean tres publicaciones compatibles del mismo acontecimiento.\par

Aquí la lectura machiana se vuelve más fuerte. La escala gravitatoria local queda seleccionada por una condición de cierre global del ciclo. El radio local de un cuanto adquiere sentido al integrarse en una periodicidad completa. La inercia, la fase y la geometría se determinan relacionalmente.\par

El alcance demostrado proporciona un mecanismo matemático concreto para una lectura machiana: una propiedad local —el acoplamiento gravitatorio— queda seleccionada por la consistencia del estado global. Las demás formulaciones históricas del principio de Mach conservan su dominio propio.\par

Después aparece la dimensión cinco.\par

El escalar \(G\) fija la intensidad gravitatoria y el portador especifica su estructura representacional. HMT construye un espacio tensorial simétrico de traza nula y dimensión cinco. Ese portador admite cinco pesos helicoidales:\par

\[
+2,\quad +1,\quad 0,\quad -1,\quad -2.
\]

Son los cinco grados de libertad naturales de un tensor simétrico sin traza en tres dimensiones antes de imponer las restricciones dinámicas de una teoría concreta.\par

Eso aclara un punto que suele confundirse. «Cinco componentes» designa el portador completo, mientras las dos helicidades designan la reducción transversal de masa nula en relatividad general.\par

Significa que el portador completo tiene cinco coordenadas irreducibles. Después, la dinámica decide qué parte de ese portador es física:\par

\begin{itemize}[leftmargin=*,itemsep=0.35em]
\item en una realización masiva de espín dos pueden sobrevivir los cinco pesos;
\item en una realización sin masa con invariancia de calibre, la proyección transversal sin traza elimina los sectores \(\pm1\) y \(0\);
\item permanecen las helicidades \(\pm2\).
\end{itemize}

Así, HMT proporciona un espacio previo a la decisión dinámica. Construye el portador completo de cinco componentes y contiene, mediante la proyección correspondiente, su reducción física a dos helicidades.\par

La belleza reside en la coexistencia ordenada de dos afirmaciones:\par

\begin{itemize}[leftmargin=*,itemsep=0.35em]
\item la geometría de espín dos tiene un portador pentacomponente;
\item la radiación gravitatoria sin masa ocupa un subespacio bidimensional.
\end{itemize}

Ambas afirmaciones forman una reducción coherente y precisa.\par

Y esa reducción conserva otra enseñanza central: el escalar \(G\) fija la intensidad común y la información de orientación vive en el portador tensorial. La intensidad reside en la constante, las helicidades residen en la representación y la genealogía completa reúne magnitud y portador.\par

